\documentclass[11pt,letterpaper]{article}
\usepackage[utf8]{inputenc}
\usepackage[T1]{fontenc}
\usepackage[spanish,es-noshorthands,es-tabla]{babel}
\usepackage[margin=2cm,top=2.5cm,bottom=2.5cm]{geometry}
\usepackage{amsmath,amssymb}
\usepackage{enumitem}
\usepackage{fancyhdr}
\usepackage{listings}
\usepackage{xcolor}
\usepackage{tcolorbox}
\usepackage{booktabs}
\usepackage{array}
\raggedbottom
\setlength{\headheight}{14pt}
\let\vspaceoriginal\vspace
\renewcommand{\vspace}[1]{\par\vspaceoriginal{#1}\par}
\pagestyle{fancy}
\fancyhf{}
\fancyhead[L]{Basic C Without Tears}
\fancyhead[C]{\NombreDocumento}
\fancyhead[R]{Stack}
\fancyfoot[L]{Basic-C-Without-Tears}
\fancyfoot[R]{\thepage}
\renewcommand{\headrulewidth}{0.4pt}
\renewcommand{\footrulewidth}{0.4pt}
\lstdefinestyle{customc}{
  language=C,
  numbers=left,
  stepnumber=1,
  numbersep=8pt,
  numberstyle=\tiny\color{gray},
  basicstyle=\footnotesize\ttfamily,
  keywordstyle=\bfseries\color{blue!80!black},
  commentstyle=\itshape\color{green!50!black},
  stringstyle=\color{red!70!black},
  breaklines=true,
  breakatwhitespace=true,
  tabsize=4,
  showstringspaces=false,
  frame=none
}
\newtcolorbox{solucionbox}{
  before=\par,
  after=\par,
  colback=white,
  colframe=black,
  arc=0mm,
  boxrule=0.6pt,
  left=3mm,
  right=3mm,
  top=2mm,
  bottom=2mm
}
\newif\ifpauta
\ifdefined\PAUTA\pautatrue\else\pautafalse\fi

\newcommand{\NombreDocumento}{Examen 02 -- Variante A}
\begin{document}
\begin{center}
    {\LARGE \textbf{Basic C Without Tears}} \\[0.2cm]
    {\Large \textbf{Examen 02 -- Variante A}} \\[0.12cm]
    \small Fundamentos de C, validación, matrices y cadenas de caracteres \\[0.08cm]
    \textbf{Autor:} Stack
\end{center}
\vspace{0.2cm}
\begin{enumerate}[label=\textbf{\arabic*.}]
\item Una estación registra temperaturas enteras durante tres días y en dos horarios:
\begin{lstlisting}[style=customc, numbers=none]
int T[3][2] = {{8, 12}, {15, 9}, {11, 14}};
\end{lstlisting}
\begin{enumerate}[label=(\alph*)]
\item Indique cuántas filas y columnas tiene la matriz. ¿Qué representan \texttt{T[1][0]} y \texttt{T[2][1]}?
\ifpauta\begin{solucionbox}La matriz tiene 3 filas y 2 columnas. \texttt{T[1][0]} vale 15 y \texttt{T[2][1]} vale 14; los índices comienzan en cero.\end{solucionbox}\else\vspace{1.5cm}\fi
\item Recorra la matriz y calcule la suma de las temperaturas de cada día (cada fila).
\ifpauta\begin{solucionbox}Las sumas por fila son: día 1, $8+12=20$; día 2, $15+9=24$; día 3, $11+14=25$.\end{solucionbox}\else\vspace{2cm}\fi
\item Escriba un fragmento en C que recorra la matriz, calcule e imprima la suma de cada fila.
\ifpauta\begin{solucionbox}\begin{lstlisting}[style=customc]
for (int i = 0; i < 3; i++) {
    int suma = 0;
    for (int j = 0; j < 2; j++) {
        suma += T[i][j];
    }
    printf("Dia %d: %d\n", i + 1, suma);
}
\end{lstlisting}El ciclo exterior selecciona la fila y el interior visita todas sus columnas.\end{solucionbox}\else\vspace{3cm}\fi
\end{enumerate}
\newpage
\item Una función recibe la edad de una persona en la variable entera \texttt{edad}. Se considera válida una edad entre 0 y 120, inclusive.
\begin{enumerate}[label=(\alph*)]
\item Escriba la condición lógica que permite comprobar que la edad es válida.
\ifpauta\begin{solucionbox}\texttt{edad >= 0 \&\& edad <= 120}. Ambas condiciones deben cumplirse.\end{solucionbox}\else\vspace{1.5cm}\fi
\item Complete un \texttt{if-else} que imprima \texttt{Edad valida} o \texttt{Edad invalida}.
\ifpauta\begin{solucionbox}\begin{lstlisting}[style=customc]
if (edad >= 0 && edad <= 120) {
    printf("Edad valida\n");
} else {
    printf("Edad invalida\n");
}
\end{lstlisting}\end{solucionbox}\else\vspace{2cm}\fi
\item Indique el resultado de la validación para \texttt{-2}, \texttt{0}, \texttt{120} y \texttt{121}.
\ifpauta\begin{solucionbox}\texttt{-2}: inválida; \texttt{0}: válida; \texttt{120}: válida; \texttt{121}: inválida. Los extremos se incluyen.\end{solucionbox}\else\vspace{2cm}\fi
\end{enumerate}
\newpage
\item Determine exactamente qué imprime el siguiente código y explique el recorrido de los ciclos.
\begin{lstlisting}[style=customc]
int M[2][3] = {{1, 2, 3}, {4, 5, 6}};
int total = 0;
for (int i = 0; i < 2; i++) {
    for (int j = 0; j < 3; j++) {
        total += M[i][j];
    }
    printf("Fila %d: %d\n", i, total);
}
\end{lstlisting}
\ifpauta\begin{solucionbox}El acumulador no se reinicia al cambiar de fila. Primera fila: suma $1+2+3=6$. Segunda fila: agrega $4+5+6$, por lo que queda 21.\par\textbf{Salida:}\\\texttt{Fila 0: 6}\\\texttt{Fila 1: 21}\end{solucionbox}\else\vspace{3cm}\fi
\item En C, una cadena termina con \texttt{'\\0'}. Escriba una función que cuente cuántas vocales minúsculas (\texttt{a,e,i,o,u}) hay en una cadena, sin modificarla.
\begin{center}\begin{tabular}{ll}\toprule\textbf{Entrada} & \textbf{Resultado}\\\midrule\texttt{"casa"} & \texttt{2}\\\texttt{"Basic C"} & \texttt{2}\\\texttt{"xyz"} & \texttt{0}\\\bottomrule\end{tabular}\end{center}
\ifpauta\begin{solucionbox}\begin{lstlisting}[style=customc]
int contar_vocales(const char texto[]) {
    int total = 0;
    for (int i = 0; texto[i] != '\0'; i++) {
        if (texto[i] == 'a' || texto[i] == 'e' ||
            texto[i] == 'i' || texto[i] == 'o' || texto[i] == 'u') {
            total++;
        }
    }
    return total;
}
\end{lstlisting}Se recorre hasta el terminador; la condición solo incluye vocales minúsculas.\end{solucionbox}\else\vspace{3cm}\fi
\end{enumerate}
\end{document}
